\section{Objet de l'annexe}
\label{annexe-reduction-objet}

Cette annexe établit le lien mathématique entre les deux premiers
tomes. Elle répond à quatre questions :

\begin{enumerate}
\item
  dans quel sens un réseau exactement homogène reste-t-il homogène ?
\item
  quel système réduit décrit alors sa dynamique ?
\item
  quels équilibres structurent le paysage autonome de ce système ?
\item
  quelle réduction scalaire est utilisée dans le Tome II pour étudier
  la résistance et l'appropriation ?
\end{enumerate}

Deux précisions sont importantes dès le départ.

Premièrement, l'homogénéité entre les nœuds du réseau et l'égalité des
composantes internes sont deux réductions différentes. Nous les
noterons respectivement \(\Delta\) et \(\mathcal D\).

Deuxièmement, la règle de priorité traite une cascade de Kairos nœud
par nœud. L'ensemble \(\Delta\) n'est donc pas invariant après chaque
saut élémentaire. Il est invariant pendant le flot et entre le début
et la fin d'une cascade complète. Cette distinction sera démontrée
explicitement.

\section{Deux réductions distinctes}
\label{annexe-reduction-deux-reductions}

\subsection{Homogénéité entre les nœuds}

Pour un réseau de \(N\) nœuds, on définit

\[
\Delta
=
\left\{
H_i=h,\;
V_i=v,\;
\tau_i=\theta,\;
P_i=p
\quad
\text{pour tout }i
\right\}.
\]

Les variables réduites

\[
h,v,\theta\in\mathbb R^d,
\qquad
p\in\mathbb R_+
\]

décrivent alors l'état commun des nœuds.

\subsection{Diagonale des composantes internes}

Lorsque \(d=2\), on introduit également

\[
\mathcal D
=
\left\{
h=(x,x),\;
v=(y,y),\;
\theta=(\vartheta,\vartheta)
\right\}.
\]

Cette seconde réduction ne concerne plus l'égalité entre les nœuds.
Elle concerne l'égalité entre les deux composantes d'un même état
réduit.

On peut donc avoir :

\[
\text{homogénéité du réseau}
\quad\text{sans}\quad
\text{égalité des composantes},
\]

ou utiliser successivement

\[
\text{réseau complet}
\longrightarrow
\Delta
\longrightarrow
\mathcal D.
\]

C'est cette double réduction qui conduit au système scalaire du
Tome II.

\section{Invariance de \texorpdfstring{$\Delta$}{Delta} pendant le flot}
\label{annexe-reduction-invariance-flot}

La matrice relationnelle normalisée est

\[
\widetilde A
=
\operatorname{Deg}(A+I)^{-1}(A+I).
\]

Elle est stochastique par lignes :

\[
\sum_{j=1}^N\widetilde A_{ij}=1.
\]

Supposons qu'à un instant \(t_0\),

\[
H_i=h,\qquad
V_i=v,\qquad
\tau_i=\theta,\qquad
P_i=p
\]

pour tout \(i\), et que l'entrée soit uniforme :

\[
u_i(t)=u(t).
\]

Alors

\[
(\widetilde A H)_i
=
\sum_j\widetilde A_{ij}h
=
h
\]

pour tout \(i\), indépendamment de la structure du graphe.

Les équations continues deviennent donc, pour chaque nœud,

\[
\dot H_i=v,
\]

\[
\dot V_i
=
-\Gamma(\theta)v
-
M(\theta)h
+
\zeta\tanh(h)
+
u(t),
\]

\[
\dot\tau_i
=
\alpha\tanh(h)-\beta\theta,
\]

\[
\dot P_i
=
\kappa_1\|h\|
+
\kappa_2\|v\|
-
\kappa_3p.
\]

Les dérivées sont les mêmes pour tous les nœuds. L'unicité de la
solution continue implique alors que les égalités entre nœuds sont
préservées.

\begin{quote}
\resultheading{Proposition}{prop:C-invariance-flot-delta}{}
Pour tout graphe, toute entrée uniforme et toute condition initiale
dans \(\Delta\), le flot reste dans \(\Delta\) jusqu'au prochain
événement.
\end{quote}

Ce résultat dépend uniquement du caractère stochastique par lignes de
\(\widetilde A\). Il ne dépend ni de la régularité du graphe ni de son
degré moyen.

\section{Effet d'une cascade de Kairos}
\label{annexe-reduction-cascade}

Supposons maintenant que l'état soit dans \(\Delta\) et que

\[
p=\Theta.
\]

Tous les nœuds atteignent donc le seuil au même instant. La règle de
priorité sélectionne d'abord un nœud \(i^\ast\). Après ce saut,

\[
P_{i^\ast}^+=r_P\Theta,
\qquad
V_{i^\ast}^+=r_Vv,
\]

tandis que les autres nœuds vérifient encore

\[
P_j=\Theta,
\qquad
V_j=v.
\]

L'état intermédiaire n'appartient généralement plus à \(\Delta\).

Le temps continu n'avance cependant pas pendant la cascade. Les
pressions des nœuds non encore traités restent égales à \(\Theta\).
Le nœud déjà traité vérifie

\[
r_P\Theta<\Theta
\]

et ne peut pas être sélectionné une seconde fois dans la même
cascade. Après au plus \(N\) sauts, chaque nœud a été traité exactement
une fois.

À la fin de la cascade,

\[
H_i^+=h,
\qquad
\tau_i^+=\theta,
\]

\[
V_i^+=r_Vv,
\qquad
P_i^+=r_P\Theta
\]

pour tout \(i\).

Le recâblage peut produire un graphe non symétrique sous l'action du
départage par les indices, mais il ne modifie pas directement les
états \(H\), \(V\), \(\tau\) et \(P\). Après la cascade,
\(\widetilde A H=h\) reste vrai puisque chaque ligne de
\(\widetilde A\) somme à un.

\begin{quote}
\resultheading{Proposition}{prop:C-invariance-cascade-delta}{}
L'ensemble \(\Delta\) est préservé par le flot et par l'application
qui associe à un état précédant une cascade l'état obtenu après la
cascade complète.
\end{quote}

Il faut donc distinguer :

\[
\boxed{
\text{invariance après chaque saut élémentaire : fausse en général}
}
\]

et

\[
\boxed{
\text{invariance entre deux cascades complètes : vraie}.
}
\]

Si l'on regroupe toute cascade instantanée en un seul saut composé,
\(\Delta\) est une variété invariante du système hybride ainsi
regroupé.

\section{Système réduit sur \texorpdfstring{$\Delta$}{Delta}}
\label{annexe-reduction-systeme-delta}

Sur \(\Delta\), la dynamique continue est

\[
\dot h=v,
\]

\[
\dot v
=
-\Gamma(\theta)v
-
M(\theta)h
+
\zeta\tanh(h)
+
u(t),
\]

\[
\dot\theta
=
\alpha\tanh(h)-\beta\theta,
\]

\[
\dot p
=
\kappa_1\|h\|
+
\kappa_2\|v\|
-
\kappa_3p.
\]

Après une cascade complète,

\[
h^+=h^-,
\qquad
\theta^+=\theta^-,
\]

\[
v^+=r_Vv^-,
\qquad
p^+=r_P\Theta.
\]

Le graphe a disparu des équations continues réduites. Il continue à
évoluer lors des Kairos, mais il ne peut pas différencier les états
tant que ceux-ci restent exactement homogènes.

La pression \(p\) ne rétroagit sur \((h,v,\theta)\) qu'au moment des
événements, par la relaxation de \(v\). Entre les événements, le
sous-système

\[
(h,v,\theta)\in\mathbb R^{3d}
\]

est autonome lorsque \(u=0\).

\section{Équilibres structurels du sous-système autonome}
\label{annexe-reduction-equilibres}

Nous appelons ici \emph{équilibre structurel} un équilibre du
sous-système \((h,v,\theta)\). La pression sera traitée séparément.

Posons

\[
k=\frac{\alpha}{\beta},
\qquad
s=\tanh(h),
\]

où \(\tanh\) est appliquée composante par composante.

À l'équilibre,

\[
v^\ast=0,
\]

et

\[
\theta^\ast
=
k\,s^\ast.
\]

L'équation de \(v\) donne

\[
M(\theta^\ast)h^\ast
=
\zeta s^\ast.
\]

Comme \(Q(\theta)=M(\theta)^{-1}\),

\[
h^\ast
=
\zeta Q(\theta^\ast)s^\ast.
\]

Or

\[
Q(\theta)
=
I_d
+
\frac{\lambda}{1+\|\theta\|^2}
\theta\theta^\top.
\]

Puisque \(\theta^\ast=ks^\ast\),

\[
\begin{aligned}
Q(\theta^\ast)s^\ast
&=
s^\ast
+
\frac{\lambda}
{1+k^2\|s^\ast\|^2}
(ks^\ast)(ks^\ast)^\top s^\ast
\\
&=
\left[
1+
\frac{\lambda k^2\|s^\ast\|^2}
{1+k^2\|s^\ast\|^2}
\right]s^\ast.
\end{aligned}
\]

Ainsi,

\[
h^\ast
=
\zeta\Lambda(h^\ast)\tanh(h^\ast),
\]

avec

\[
\Lambda(h)
=
1+
\frac{\lambda k^2\|\tanh(h)\|^2}
{1+k^2\|\tanh(h)\|^2}.
\]

Le facteur \(\zeta\Lambda(h^\ast)\) est commun à toutes les
composantes.

\section{Pourquoi toutes les composantes non nulles ont la même
amplitude}
\label{annexe-reduction-amplitude-commune}

Une composante non nulle \(y\) d'un équilibre satisfait

\[
y=C\tanh(y),
\qquad
C>0.
\]

Pour \(y>0\), cette équation équivaut à

\[
C=\frac{y}{\tanh y}.
\]

La fonction

\[
\varphi(y)=\frac{y}{\tanh y}
\]

est strictement croissante sur \((0,\infty)\). En effet,

\[
\varphi'(y)
=
\frac{\tanh y-y\,\operatorname{sech}^2y}
{\tanh^2y}.
\]

Posons

\[
g(y)
=
\tanh y-y\,\operatorname{sech}^2y.
\]

On a \(g(0)=0\) et

\[
g'(y)
=
2y\,\operatorname{sech}^2y\,\tanh y
>0
\]

pour \(y>0\). Donc \(g(y)>0\), puis \(\varphi'(y)>0\).

Pour une valeur donnée de \(C\), il existe donc au plus une solution
positive. Par imparité de \(\tanh\), les solutions non nulles sont
\(\pm x\).

Si \(q\) composantes de \(h^\ast\) sont non nulles, alors

\[
\|s^\ast\|^2
=
q\tanh^2x.
\]

L'amplitude commune \(x>0\) vérifie

\[
\boxed{
x
=
\zeta\Lambda_q(x)\tanh x
}
\]

avec

\[
\boxed{
\Lambda_q(x)
=
1+
\frac{\lambda qk^2\tanh^2x}
{1+qk^2\tanh^2x}.
}
\]

\begin{quote}
\resultheading{Proposition}{prop:C-equilibres-structurels}{}
Tout équilibre structurel vérifie
\[
v^\ast=0,
\qquad
\theta^\ast=k\tanh(h^\ast),
\]
et les composantes de \(h^\ast\) appartiennent à
\[
\{0,\pm x_q\},
\]
où \(x_q\) est une racine positive de l'équation précédente et \(q\)
est le nombre de composantes non nulles.
\end{quote}

Cette proposition classe la forme possible des équilibres. Le nombre
de racines positives de l'équation scalaire dépend des paramètres.

\section{Développement près de l'origine}
\label{annexe-reduction-developpement-origine}

Pour \(x\) petit,

\[
\tanh x
=
x-\frac{x^3}{3}+O(x^5),
\]

et

\[
\tanh^2x
=
x^2+O(x^4).
\]

Par conséquent,

\[
\Lambda_q(x)
=
1+\lambda qk^2x^2+O(x^4).
\]

L'équation non nulle, divisée par \(x\), devient

\[
1
=
\zeta
\left[
1+
\left(
\lambda qk^2-\frac13
\right)x^2
+
O(x^4)
\right].
\]

Lorsque

\[
\lambda qk^2>\frac13,
\]

la petite racine est donc approximativement

\[
\boxed{
x_{u,q}
\approx
\sqrt{
\frac{1/\zeta-1}
{\lambda qk^2-1/3}
}.
}
\]

Cette formule est une approximation locale. Elle ne démontre pas à
elle seule l'existence d'une racine haute.

\section{Condition simple d'extinction structurelle}
\label{annexe-reduction-extinction-structurelle}

Pour tout \(x>0\),

\[
\tanh x<x,
\]

et

\[
\Lambda_q(x)<1+\lambda.
\]

Si

\[
\zeta(1+\lambda)\leq1,
\]

alors

\[
\zeta\Lambda_q(x)\tanh x
<
\zeta(1+\lambda)x
\leq x.
\]

L'équation

\[
x=\zeta\Lambda_q(x)\tanh x
\]

ne possède donc aucune solution positive.

\begin{quote}
\resultheading{Proposition}{prop:C-extinction-structurelle}{}
Si
\[
\zeta(1+\lambda)\leq1,
\]
l'origine est le seul équilibre structurel sur \(\Delta\).
\end{quote}

La condition est suffisante. Lorsque
\(\zeta(1+\lambda)>1\), l'existence d'un équilibre actif dépend encore
de la forme complète de l'équation.

\section{Stabilité locale de l'origine}
\label{annexe-reduction-stabilite-origine}

À l'origine,

\[
M(0)=I_d,
\qquad
\Gamma(0)=\gamma_0,
\qquad
D\tanh(0)=I_d.
\]

Le terme

\[
\|\theta\|v
\]

est d'ordre deux au voisinage de \((\theta,v)=(0,0)\). Il ne contribue
donc pas à la linéarisation, malgré la non-différentiabilité de la
norme en zéro.

La linéarisation est

\[
\delta\dot h=\delta v,
\]

\[
\delta\dot v
=
-(1-\zeta)\delta h
-
\gamma_0\delta v,
\]

\[
\delta\dot\theta
=
\alpha\delta h
-
\beta\delta\theta.
\]

La matrice est triangulaire par blocs. Son polynôme caractéristique
est

\[
\boxed{
(s+\beta)^d
\left(
s^2+\gamma_0s+1-\zeta
\right)^d.
}
\]

Comme \(\beta>0\) et \(\gamma_0>0\), l'origine est localement
exponentiellement stable si

\[
\boxed{\zeta<1}.
\]

Elle est instable si \(\zeta>1\). Le cas \(\zeta=1\) est
non hyperbolique.

\section{Seuil de trace qui déstabilise le repos}
\label{annexe-reduction-seuil-trace}

Le Tome II considère un système maintenu près de

\[
h=0,
\qquad
v=0,
\]

alors que sa trace peut être chargée par une influence extérieure.

Dans la direction de la trace, la valeur propre du préconditionneur
est

\[
m(r)
=
\frac{1+r^2}
{1+(1+\lambda)r^2},
\qquad
r=\|\theta\|.
\]

Pour une trace momentanément fixée, la dynamique linéaire de la
position dans cette direction est

\[
\ddot x
+
\Gamma(r)\dot x
+
\bigl(m(r)-\zeta\bigr)x
=
0.
\]

Le repos perd sa stabilité lorsque

\[
m(r)=\zeta.
\]

L'équation donne

\[
1+r^2
=
\zeta
\left[
1+(1+\lambda)r^2
\right].
\]

Ainsi,

\[
r^2
=
\frac{1-\zeta}
{\zeta(1+\lambda)-1}.
\]

Un seuil fini positif existe si et seulement si

\[
\frac{1}{1+\lambda}<\zeta<1.
\]

On obtient alors

\[
\boxed{
\tau_{\mathrm{ign}}
=
\sqrt{
\frac{1-\zeta}
{\zeta(1+\lambda)-1}
}.
}
\]

Pour les paramètres de référence,

\[
\zeta=0{,}8,
\qquad
\lambda=2,
\]

donc

\[
\boxed{
\tau_{\mathrm{ign}}
=
0{,}377964\ldots
}
\]

Le seuil concerne la norme de la trace. Sur la diagonale interne
\(\theta=(\vartheta,\vartheta)\),

\[
\|\theta\|=\sqrt2\,|\vartheta|.
\]

\section{Extinction globale lorsque la trace n'agit pas sur le rappel}
\label{annexe-reduction-extinction-lambda-zero}

Prenons

\[
\lambda=0.
\]

Alors

\[
M(\theta)=I_d
\]

quelle que soit la trace. Le sous-système \((h,v)\) vérifie

\[
\dot h=v,
\]

\[
\dot v
=
-\Gamma(\theta)v
-h
+
\zeta\tanh h.
\]

Supposons \(\zeta<1\). Définissons

\[
\mathcal V(h,v)
=
\frac12\|v\|^2
+
\frac12\|h\|^2
-
\zeta
\sum_{k=1}^d\log\cosh(h_k).
\]

Comme

\[
\frac{\partial}{\partial h_k}
\log\cosh(h_k)
=
\tanh(h_k),
\]

sa dérivée est

\[
\begin{aligned}
\dot{\mathcal V}
&=
v\cdot\dot v
+
\bigl(h-\zeta\tanh h\bigr)\cdot\dot h
\\
&=
-\Gamma(\theta)\|v\|^2
\\
&\leq
-\gamma_0\|v\|^2.
\end{aligned}
\]

La fonction \(\mathcal V\) est coercive. En effet,

\[
\log\cosh y\leq |y|,
\]

et le terme quadratique domine lorsque \(\|h\|\to\infty\).

Elle est également strictement convexe en \(h\), car

\[
D_h^2\mathcal V
=
I_d
-
\zeta\,
\operatorname{diag}
\bigl(\operatorname{sech}^2h_k\bigr)
\succeq
(1-\zeta)I_d.
\]

Son unique minimum est donc \(h=0\), \(v=0\).

L'ensemble où \(\dot{\mathcal V}=0\) impose \(v=0\). Pour qu'une
trajectoire reste dans cet ensemble, il faut aussi

\[
\dot v
=
-h+\zeta\tanh h
=
0.
\]

Comme \(\zeta<1\), cette équation n'a que la solution \(h=0\). Enfin,

\[
\dot\theta=-\beta\theta
\]

sur cet ensemble.

Le principe d'invariance de LaSalle donne donc

\[
\boxed{
(h(t),v(t),\theta(t))
\longrightarrow
(0,0,0).
}
\]

L'état actif du régime de référence n'existe donc pas sans
l'affaiblissement du rappel produit par la trace.

\section{Stabilité d'un équilibre en dimension un}
\label{annexe-reduction-stabilite-d1}

En dimension \(d=1\), écrivons

\[
x=h,
\qquad
\vartheta=\theta.
\]

La dynamique autonome est

\[
\dot x=v,
\]

\[
\dot v
=
-\Gamma(\vartheta)v
-
m(\vartheta)x
+
\zeta\tanh x,
\]

\[
\dot\vartheta
=
\alpha\tanh x-\beta\vartheta,
\]

avec

\[
m(\vartheta)
=
\frac{1+\vartheta^2}
{1+(1+\lambda)\vartheta^2}.
\]

À un équilibre non nul,

\[
\vartheta^\ast
=
k\tanh x^\ast.
\]

Posons

\[
s_\ast
=
\operatorname{sech}^2x^\ast,
\]

\[
c_\ast
=
\zeta s_\ast,
\]

\[
m_\ast
=
m(\vartheta^\ast),
\]

et

\[
b_\ast
=
-m'(\vartheta^\ast)x^\ast
=
\frac{
2\lambda\vartheta^\ast x^\ast
}{
\left[
1+(1+\lambda)(\vartheta^\ast)^2
\right]^2
}.
\]

Le produit \(\vartheta^\ast x^\ast\) est positif, donc \(b_\ast>0\).

Le Jacobien est

\[
J_\ast
=
\begin{pmatrix}
0 & 1 & 0\\
-(m_\ast-c_\ast) & -\Gamma_\ast & b_\ast\\
\alpha s_\ast & 0 & -\beta
\end{pmatrix},
\]

où

\[
\Gamma_\ast
=
\gamma_0+\gamma_1|\vartheta^\ast|.
\]

Son polynôme caractéristique est

\[
p(s)
=
s^3+a_2s^2+a_1s+a_0,
\]

avec

\[
a_2
=
\Gamma_\ast+\beta,
\]

\[
a_1
=
\Gamma_\ast\beta
+
m_\ast-c_\ast,
\]

\[
a_0
=
\beta(m_\ast-c_\ast)
-
\alpha b_\ast s_\ast.
\]

Pour un polynôme cubique réel, les conditions de Routh--Hurwitz sont

\[
a_2>0,
\qquad
a_1>0,
\qquad
a_0>0,
\qquad
a_2a_1>a_0.
\]

\subsection{Lien entre \texorpdfstring{$a_0$}{a0} et la branche d'équilibres}

Définissons

\[
G(x)
=
m\!\left(k\tanh x\right)x
-
\zeta\tanh x.
\]

Les équilibres non nuls sont les zéros de \(G\). Par dérivation,

\[
\begin{aligned}
G'(x)
={}&
m_\ast
-
\zeta s_\ast
+
m'(\vartheta^\ast)
\left(
ks_\ast
\right)x^\ast
\\
={}&
m_\ast-c_\ast
-
\frac{\alpha}{\beta}
b_\ast s_\ast.
\end{aligned}
\]

Par conséquent,

\[
\boxed{
a_0
=
\beta G'(x^\ast).
}
\]

Si la branche traverse l'axe avec \(G'(x^\ast)<0\), alors
\(a_0<0\) et l'équilibre est instable.

Pour les paramètres de référence, la racine basse possède
\(a_0<0\), tandis que la racine haute satisfait toutes les conditions
de Routh--Hurwitz.

\section{Classification en dimension deux}
\label{annexe-reduction-classification-d2}

Pour les paramètres de référence,

\[
\zeta=0{,}8,
\qquad
\lambda=2,
\qquad
\alpha=0{,}45,
\qquad
\beta=0{,}12,
\]

\[
\gamma_0=0{,}25,
\qquad
\gamma_1=1{,}25,
\qquad
d=2.
\]

L'équation scalaire possède les racines positives suivantes :

\[
\begin{array}{c|cc}
q & \text{racine basse} & \text{racine haute}\\
\hline
1 & 0{,}1021426 & 2{,}2376194\\
2 & 0{,}0717413 & 2{,}2958568
\end{array}
\]

Le statut exact de ce tableau est numérique : les racines sont
obtenues par isolation puis résolution de l'équation scalaire. Leur
forme possible, en revanche, est donnée analytiquement par la
Proposition~\ref{prop:C-equilibres-structurels}.

On obtient :

\begin{itemize}
\item
  l'origine ;
\item
  pour \(q=1\), des points d'axe
  \[
  (\pm x,0),
  \qquad
  (0,\pm x);
  \]
\item
  pour \(q=2\), des points à deux composantes actives
  \[
  (\pm x,\pm x),
  \]
  avec choix indépendant des signes.
\end{itemize}

Chaque famille existe avec une racine basse et une racine haute.

\section{Pourquoi les points d'axe ne sont pas attractifs}
\label{annexe-reduction-instabilite-axe}

Considérons un point d'axe

\[
h^\ast=(x^\ast,0),
\qquad
\theta^\ast=(\vartheta^\ast,0),
\]

avec

\[
\vartheta^\ast=k\tanh x^\ast.
\]

La dynamique longitudinale, dans la première composante, est celle du
système scalaire précédent.

Étudions une perturbation dans la seconde composante. À l'équilibre,

\[
M(\theta^\ast)
=
\begin{pmatrix}
m(\vartheta^\ast)&0\\
0&1
\end{pmatrix}.
\]

La rotation infinitésimale de la direction de la trace produit le
coefficient

\[
b_\perp
=
\frac{
\lambda\vartheta^\ast x^\ast
}{
1+(1+\lambda)(\vartheta^\ast)^2
}.
\]

Le bloc transverse est

\[
J_\perp^{\mathrm{axe}}
=
\begin{pmatrix}
0&1&0\\
-(1-\zeta)&-\Gamma_\ast&b_\perp\\
\alpha&0&-\beta
\end{pmatrix}.
\]

Son terme constant est

\[
\boxed{
a_{0,\perp}^{\mathrm{axe}}
=
\beta(1-\zeta)-\alpha b_\perp.
}
\]

Pour la racine haute \(x^\ast=2{,}2376194\),

\[
\beta(1-\zeta)=0{,}024,
\]

tandis que

\[
\alpha b_\perp
=
0{,}178700\ldots
\]

Ainsi,

\[
a_{0,\perp}^{\mathrm{axe}}
=
-0{,}154700\ldots<0.
\]

Le bloc possède une valeur propre positive :

\[
\max\Re\operatorname{spec}
\left(
J_\perp^{\mathrm{axe}}
\right)
=
0{,}113981\ldots
\]

Le point d'axe haut est donc stable le long de l'axe, mais instable
vis-à-vis de l'allumage de la composante éteinte.

\section{Stabilité des points à deux composantes actives}
\label{annexe-reduction-stabilite-diagonale}

Considérons, par symétrie,

\[
h^\ast=(x^\ast,x^\ast),
\qquad
\theta^\ast=(\vartheta^\ast,\vartheta^\ast).
\]

Les autres choix de signes ont le même spectre.

Posons

\[
e_+
=
\frac{1}{\sqrt2}(1,1),
\qquad
e_-
=
\frac{1}{\sqrt2}(1,-1).
\]

La linéarisation se décompose en deux blocs de dimension trois :

\begin{itemize}
\item
  un bloc longitudinal dans la direction \(e_+\) ;
\item
  un bloc de désaccord interne entre composantes, dans la direction \(e_-\).
\end{itemize}

Notons

\[
r_\ast
=
\|\theta^\ast\|
=
\sqrt2\,|\vartheta^\ast|,
\]

\[
h_\ast^\parallel
=
\sqrt2\,|x^\ast|,
\]

\[
s_\ast
=
\operatorname{sech}^2x^\ast.
\]

\[
\Gamma_\ast
=
\gamma_0+\gamma_1 r_\ast .
\]

Le symbole \(\Gamma_\ast\) est ici évalué en \(\|\theta^\ast\|=r_\ast\), et
non en \(|\vartheta^\ast|\) comme dans la section en dimension un.

Pour le bloc longitudinal,

\[
\kappa_+
=
m(r_\ast)-\zeta s_\ast,
\]

et

\[
b_+
=
\frac{
2\lambda r_\ast h_\ast^\parallel
}{
\left[
1+(1+\lambda)r_\ast^2
\right]^2
}.
\]

Pour le bloc de désaccord interne,

\[
\kappa_-
=
1-\zeta s_\ast,
\]

et

\[
b_-
=
\frac{
\lambda r_\ast h_\ast^\parallel
}{
1+(1+\lambda)r_\ast^2
}.
\]

Les deux blocs ont la forme

\[
J_\pm
=
\begin{pmatrix}
0&1&0\\
-\kappa_\pm&-\Gamma_\ast&b_\pm\\
\alpha s_\ast&0&-\beta
\end{pmatrix}.
\]

Leurs polynômes caractéristiques sont

\[
p_\pm(s)
=
s^3
+
(\Gamma_\ast+\beta)s^2
+
(\Gamma_\ast\beta+\kappa_\pm)s
+
\left(
\beta\kappa_\pm
-
\alpha b_\pm s_\ast
\right).
\]

Pour la racine basse

\[
x^\ast=0{,}0717413,
\]

le bloc longitudinal possède un terme constant négatif :

\[
a_{0,+}
=
-0{,}0332873\ldots
\]

et l'équilibre est instable.

Pour la racine haute

\[
x^\ast=2{,}2958568,
\]

les coefficients sont :

\[
\begin{array}{c|ccc}
 & a_2 & a_1 & a_0\\
\hline
\text{longitudinal}
&
6{,}866096
&
1{,}119207
&
0{,}036982
\\
\text{désaccord}
&
6{,}866096
&
1{,}777746
&
0{,}108829
\end{array}
\]

Les quatre conditions de Routh--Hurwitz sont satisfaites dans les deux
blocs. Une évaluation séparée donne
\[
\max\Re\operatorname{spec}(J_+)
=
-0{,}0458575\ldots,
\qquad
\max\Re\operatorname{spec}(J_-)
=
-0{,}0971661\ldots.
\]
Le spectre complet vérifie donc

\[
\boxed{
\max\Re\operatorname{spec}(J_\ast)
=
-0{,}0458575\ldots
}
\]

La racine haute à deux composantes actives est donc localement
exponentiellement stable sur le système réduit.

\section{Paysage autonome aux paramètres de référence}
\label{annexe-reduction-paysage-reference}

Les calculs précédents donnent le paysage suivant pour
\((h,v,\theta)\) sur \(\Delta\).

\[
\begin{array}{l|c|c}
\text{famille} & \|h^\ast\| & \max\Re\operatorname{spec}(J_\ast)\\
\hline
\text{origine}
&
0
&
-0{,}12
\\
\text{axe, racine basse}
&
0{,}102143
&
+0{,}141613
\\
\text{axe, racine haute}
&
2{,}237619
&
+0{,}113981
\\
\text{deux composantes, racine basse}
&
0{,}101458
&
+0{,}143662
\\
\text{deux composantes, racine haute}
&
3{,}246832
&
-0{,}0458575
\end{array}
\]

Il existe donc simultanément :

\begin{itemize}
\item
  un repos localement stable ;
\item
  quatre états actifs stables, correspondant aux quatre choix de signes
  \[
  (\pm x_2,\pm x_2);
  \]
\item
  des équilibres instables qui constituent des candidats naturels pour structurer localement les séparations entre bassins.
\end{itemize}

La coexistence de plusieurs attracteurs locaux implique que leurs
bassins contiennent des ouverts disjoints. L'état atteint dépend donc
de la condition initiale.

Cette conclusion est analytique quant à la coexistence et à la
stabilité locale, modulo l'évaluation numérique des racines et des
inégalités de Routh--Hurwitz. La géométrie globale des séparatrices
n'est pas démontrée.

\section{La pression et le sens du mot ``puits''}
\label{annexe-reduction-pression-puits}

L'origine est un véritable équilibre du système hybride complet :

\[
h=v=\theta=0,
\qquad
p=0.
\]

Pour un état actif structurel,

\[
h=h^\ast,
\qquad
v=0,
\qquad
\theta=\theta^\ast,
\]

la pression vérifie

\[
\dot p
=
\kappa_1\|h^\ast\|
-
\kappa_3p.
\]

Sa valeur d'équilibre continue serait

\[
p_\infty
=
\frac{\kappa_1}{\kappa_3}
\|h^\ast\|.
\]

Pour le puits actif de référence,

\[
\|h^\ast\|
=
3{,}246832,
\]

donc

\[
p_\infty
=
3{,}896198
>
\Theta=1{,}25.
\]

La pression atteint donc périodiquement le seuil. Comme \(v=0\), la
relaxation

\[
v^+=r_Vv^-
\]

ne modifie pas l'état structurel. Seules la pression et la structure
du graphe évoluent.

Après une cascade,

\[
p(0)=r_P\Theta.
\]

La solution est

\[
p(t)
=
p_\infty
+
\left(
r_P\Theta-p_\infty
\right)e^{-\kappa_3t}.
\]

Le prochain seuil est atteint après

\[
\boxed{
\Delta_p
=
\frac{1}{\kappa_3}
\log
\left(
\frac{p_\infty-r_P\Theta}
{p_\infty-\Theta}
\right).
}
\]

Pour les paramètres de référence,

\[
\Delta_p
=
0{,}535537\ldots
\]

Ainsi, le ``puits actif'' est exactement un point fixe dans la
projection \((h,v,\theta)\), mais il correspond dans le système
hybride complet à un régime invariant comportant des cascades
répétées de Kairos.

La définition de l'appropriation par les bassins reste correcte si
l'attracteur autonome \(\mathcal A\) désigne cet ensemble hybride.
Par commodité, le corps du texte utilise souvent le mot ``puits'' pour
sa projection structurale.

\section{Réduction scalaire utilisée dans le Tome II}
\label{annexe-reduction-scalaire-tome-ii}

Considérons maintenant la diagonale interne

\[
h=(x,x),
\qquad
v=(y,y),
\qquad
\theta=(\vartheta,\vartheta).
\]

La fonction \(\tanh\) préserve cette diagonale :

\[
\tanh(h)
=
(\tanh x,\tanh x).
\]

La trace est parallèle à \(h\), et

\[
\|\theta\|
=
\sqrt2\,|\vartheta|.
\]

Dans la direction de la trace,

\[
M(\theta)h
=
m(\sqrt2|\vartheta|)h,
\]

avec

\[
m(r)
=
\frac{1+r^2}
{1+(1+\lambda)r^2}.
\]

La dynamique autonome devient

\[
\dot x=y,
\]

\[
\dot y
=
-\Gamma(\sqrt2|\vartheta|)y
-
m(\sqrt2|\vartheta|)x
+
\zeta\tanh x,
\]

\[
\dot\vartheta
=
\alpha\tanh x-\beta\vartheta.
\]

Si un système extérieur écrit dans la trace avec une projection
constante \(w\) et une intensité \(g\), seule la dernière équation est
modifiée :

\[
\boxed{
\dot\vartheta
=
\alpha\tanh x
+
gw
-
\beta\vartheta.
}
\]

C'est le système scalaire utilisé pour étudier la branche résistante.

La diagonale \(\mathcal D\) est invariante parce que :

\begin{itemize}
\item
  les deux composantes reçoivent la même entrée ;
\item
  \(\tanh\) agit composante par composante ;
\item
  \(M(\theta)\) conserve la direction de \(\theta\) ;
\item
  les conditions initiales ont des composantes égales.
\end{itemize}

La réduction est exacte sur \(\mathcal D\). Elle ne décrit pas les
perturbations qui rendent les deux composantes différentes.

\section{Portée pour le réseau complet}
\label{annexe-reduction-portee-reseau}

Pour un graphe figé, la stabilité transverse d'un état homogène peut
être étudiée en décomposant les perturbations selon les modes propres
de \(\widetilde A\).

Lorsque le graphe est non orienté,
\(\widetilde A\) est semblable à la matrice symétrique

\[
\operatorname{Deg}(A+I)^{-1/2}
(A+I)
\operatorname{Deg}(A+I)^{-1/2}.
\]

Ses valeurs propres sont donc réelles. Soit \(\rho\) l'une d'elles.

Autour d'un état homogène \((h^\ast,0,\theta^\ast)\), le bloc de flot
associé au mode \(\rho\) est

\[
J(\rho)
=
\begin{pmatrix}
0&I_d&0\\
-M_\ast+\zeta\rho S_\ast
&
-\Gamma_\ast I_d
&
-D_\ast
\\
\alpha\rho S_\ast
&
0
&
-\beta I_d
\end{pmatrix},
\]

où

\[
S_\ast
=
\operatorname{diag}
\left(
\operatorname{sech}^2h_k^\ast
\right),
\]

et \(D_\ast\) est l'application linéaire

\[
\eta
\longmapsto
DM(\theta^\ast)[\eta]h^\ast.
\]

Le mode \(\rho=1\) est le mode homogène et redonne le Jacobien étudié
plus haut. Les blocs \(J(\rho)\) avec \(\rho\neq1\) sont les
\emph{blocs modaux transverses au réseau} : ils comparent les nœuds du réseau.
Ils ne doivent pas être confondus avec le bloc de désaccord interne
\(J_-\), qui compare les deux composantes internes d'un même état réduit.

Cette formule démontre la réduction modale du flot pour un graphe
figé. Elle ne constitue pas, à elle seule, une preuve de stabilité du
système hybride avec recâblage : une analyse complète devrait inclure
la sensibilité des instants de Kairos et les matrices de saltation.

Aux paramètres et graphes testés, les modes transverses décroissent
numériquement. Ce fait appartient au registre des validations
numériques, non à celui des théorèmes établis ici.

\section{Résultats établis et limites}
\label{annexe-reduction-resultats-limites}

Cette annexe établit analytiquement :

\begin{itemize}
\item
  l'invariance de \(\Delta\) pendant le flot ;
\item
  son retour après chaque cascade complète ;
\item
  le système réduit sur \(\Delta\) ;
\item
  la classification structurelle des équilibres ;
\item
  la stabilité locale de l'origine ;
\item
  le seuil
  \(\tau_{\mathrm{ign}}\) ;
\item
  l'extinction globale lorsque \(\lambda=0\) et \(\zeta<1\) ;
\item
  les Jacobiennes et les critères de stabilité des branches ;
\item
  la réduction scalaire exacte utilisée dans le Tome II.
\end{itemize}

Les valeurs numériques des racines et des spectres ont le statut
\texttt{[PROP-NUM]}.

Ne sont pas démontrés ici :

\begin{itemize}
\item
  la géométrie globale des bassins ;
\item
  la stabilité hybride transverse avec instants d'événements
  variables et graphe évolutif ;
\item
  l'existence et la stabilité des états mixtes du réseau complet ;
\item
  une classification globale du paysage pour tous les paramètres.
\end{itemize}

Le résultat nécessaire à la suite est néanmoins acquis :

\[
\boxed{
\begin{minipage}{0.86\linewidth}
Le système autonome réduit possède simultanément un repos stable et quatre régimes actifs stables. Les bassins de ces attracteurs contiennent des ouverts disjoints ; la géométrie globale de leurs séparations n'est pas établie ici. Une influence extérieure peut modifier l'accessibilité de ces régimes sans les créer artificiellement : après découplage, le régime atteint appartient encore au paysage autonome du système receveur.
\end{minipage}
}
\]
